\documentclass[aps,prl,reprint]{revtex4-2}
% Explicit, opt-in helpers for local typography previews.
% Choose the helper only after journal AND article type are known.
% Heading presence and wording remain author decisions.
% Pass heading words without terminal punctuation, except Note added. below.
% Never redefine \section globally with these helpers.
\newcommand{\prlrunin}[1]{%
  \par\noindent\hspace*{\parindent}%
  {\normalfont\itshape #1---}\ignorespaces}
\newcommand{\apsletterrunin}[1]{%
  \par\noindent\hspace*{\parindent}%
  {\normalfont\itshape #1.}\ \ignorespaces}
\newcommand{\apssecondaryrunin}[1]{%
  \par\noindent\hspace*{\parindent}%
  {\normalfont #1:}\quad\ignorespaces}
\newcommand{\apsdeepprlrunin}[1]{%
  \par\noindent\hspace*{\parindent}%
  {\normalfont\itshape #1.}\ \ignorespaces}

\begin{document}
\title{Synthetic layout examples: PRL with an unheaded opening}
\author{Typography demonstration only}\noaffiliation
\maketitle
This synthetic opening is intentionally unheaded. It demonstrates that selecting PRL typography does not require an Introduction label. The paragraph reports no scientific result and is not a reproduction of a published article.

The paper may continue naturally for as many paragraphs as the actual argument requires. Nothing in the heading helpers inserts a missing section or dictates an IMRaD outline.

\prlrunin{Illustrative setup}
Only an author-selected descriptive heading is printed here. It is italic, indented, unnumbered, and followed by an italic em dash without a preceding period.

\prlrunin{Outlook}
This optional closing heading replaces a Conclusion label in this example. The body text resumes upright on the same baseline. The wording remains an editorial choice for the actual manuscript.

\begin{acknowledgments}
\prlrunin{Acknowledgments}
Illustrative placeholder only. No real contribution or support is asserted.
\end{acknowledgments}
\prlrunin{Data availability}
Illustrative placeholder only. A real manuscript needs a truthful statement based on its actual research.
\end{document}
